\section{Kummer calculations with uniform supports}\label{sec:kummer}

To compute supports on the extension-space quotient, we need comparison
maps that are uniform as analytic bounds and compact parameters vary.
The first goal is a bounded Kummer homotopy on nested annuli. Disc,
affine-line, and shrinking-graph calculations will then give the exact
support and neighborhood operations used for rank deletion.

This section concerns the \emph{untilted} affine line.  Put
$C=\Cp$ and let $S=\Spa(R,R^+)$ be an affinoid perfectoid space over
$\Cflat$, with its specified untilt over $C$.  Write
\[
 \mathbb A_S=(\mathbb A^1_C)^\diamond\times_{\operatorname{Spd}C}S.
\]
The sheaf $\cO^\flat$ on this space assigns to a perfectoid test object
its characteristic-$p$ structural ring.  In particular, it is not the
structural sheaf on the perfected characteristic-$p$ affine line.
The latter appears in \cref{sec:laurent} and has a different support
calculation.  The identification of the tilted sheaf on perfectoid
charts follows from \cite[Definition~5.10 and Lemma~5.11]{scholzeHodge2013};
we use its v-acyclicity in the form
\cite[Proposition~8.8]{scholzeDiamonds2017}.
The toric Kummer antecedent is
\cite[Example~4.4, Lemma~4.5(i), and the proof of Lemma~5.6]{scholzeHodge2013},
with the countable-cover convention of
\cite[item~(1)]{scholzeHodgeCorrigendum2016}.  We adapt that construction
to relative annuli and prove the uniform support estimates below.

All radii below belong to $|C^\times|$.  Closed inequalities are
interleaved with slightly larger open inequalities.  Thus, for a
relative bound $K$ in an open space $X$, the notation
\begin{equation}\label{eq:kum-support-fiber}
 \RGamma_K(X,\mathscr F)
 =\operatorname{fib}\bigl(\RGamma(X,\mathscr F)
           \longrightarrow\RGamma(X\setminus K,\mathscr F)\bigr)
\end{equation}
can be computed using separated inner and outer bounds.  Compact
support means the filtered colimit of these fibers over relatively
compact bounds.  For a relative calculation the bounds are in the
affine-line coordinates; the base and all compact parameters are
retained.  These conventions avoid imposing a particular sharp
boundary at higher-rank points.

\subsection{Complete coefficient rings and inverse limits}

For a compact profinite set $\mathcal D$, the ring
\[
 R_{\mathcal D}=C(\mathcal D,\Cflat)
\]
is equipped with the supremum norm.  More generally we may replace
$R$ by $C(\mathcal D,R)$.  This is again a complete uniform perfectoid
algebra: Frobenius is bijective, its inverse is continuous, and the
power-bounded elements are the continuous functions into $R^\circ$.
All Laurent norms used below are supremum norms with coefficients in
this ring.  In particular, a bound for a scalar multiplication
operator is automatically uniform in $\mathcal D$.

The dense-tower criterion below is a special case of
\citet[Proposition~3.5]{lebras2018}. We include the argument to
retain uniformity in compact parameters and the degreewise statement
for complexes.

\begin{lemma}[Countable Banach inverse limits]\label[lemma]{lem:kum-banach-limits}
Let $(E_j,f_j)$ be a countable inverse system of complete
nonarchimedean normed vector spaces, with continuous linear maps
$f_j:E_{j+1}\to E_j$ having dense image.  Then
$R^1\!\varprojlim_j E_j=0$.  The assertion holds after replacing
$E_j$ by $C(\mathcal D,E_j)$ with its supremum norm for any compact
profinite $\mathcal D$.  Consequently a countable inverse system of
complexes having this property in each degree is computed by its
termwise inverse limit.
\end{lemma}

\begin{proof}
The countable derived limit is represented by
\[
 \prod_{j\geq0}E_j\longrightarrow\prod_{j\geq0}E_j,
 \qquad (v_j)_j\longmapsto(v_j-f_j(v_{j+1}))_j.
\]
We prove surjectivity.  Given $(a_j)_j$, construct finite strings
$v^{(h)}_0,\ldots,v^{(h)}_h$ satisfying
$v^{(h)}_j-f_j(v^{(h)}_{j+1})=a_j$ for $j<h$.
Having constructed such a string, use density to choose
$v^{(h+1)}_{h+1}$ so that
\[
 e_h=a_h+f_h(v^{(h+1)}_{h+1})-v^{(h)}_h
\]
is as small as desired.  Define the earlier terms by the required
relations.  Their changes are
$f_j\cdots f_{h-1}(e_h)$ for $j<h$ and $e_h$ for $j=h$.
By continuity of these finitely many maps, choose $e_h$ so small that
all those changes have norm at most $2^{-h}$.  For each fixed $j$ the
sequence $v^{(h)}_j$ is Cauchy.  Completeness and continuity give limits
$v_j$ satisfying every required relation.

For the assertion with parameters, a continuous map from
$\mathcal D$ to $E_j$ can be approximated uniformly by a map constant
on a finite clopen partition.  Approximate its finitely many values
by images from $E_{j+1}$.  Thus
$C(\mathcal D,E_{j+1})\to C(\mathcal D,E_j)$ has dense image, and the
same proof applies.  Finally the displayed product complex computes
the derived inverse limit degreewise.  Surjectivity in each degree
identifies it with the kernel complex, which is the termwise inverse
limit.  No closed-image assumption is needed.
\end{proof}

We will also use the elementary countable Milnor exact sequence
\begin{equation}\label{eq:kum-milnor}
 0\longrightarrow R^1\!\varprojlim_j H^{i-1}(E_j^\bullet)
 \longrightarrow H^i(\Rlim_j E_j^\bullet)
 \longrightarrow \varprojlim_j H^i(E_j^\bullet)
 \longrightarrow0.
\end{equation}
In particular, a pro-zero cohomology system contributes neither a
limit nor a first derived limit.  Here pro-zero means that for each
target index a sufficiently late transition to that index is zero.

\subsection{A uniform annular homotopy}

Choose compatible primitive roots of unity
$(1,\zeta_p,\zeta_{p^2},\ldots)$ and let
$\epsilon\in\cO_{\Cflat}^{\times}$ be the corresponding tilted
element.  Thus $\epsilon^\sharp=1$, $\epsilon\ne1$, and
\[
 \kappa=-\log|\epsilon-1|>0.
\]
A choice of this sequence identifies $\Zp(1)$ with $\Zp$; it also
chooses all the geometric orientation generators in this section.
We make no arithmetic trivialization assertion.

For $0<a<b$, let $A_S[a,b]\subset\mathbb A_S$ be the relative
annulus $a\leq|T|\leq b$.  Define the complete Laurent algebra
\begin{equation}\label{eq:kum-annular-algebra}
 M_R[a,b]=\left\{
  \sum_{u\in\Z[1/p]}c_uT^u:
  \|c_u\|\max(a^u,b^u)\longrightarrow0
 \right\}.
\end{equation}
The arrow to zero means that only finitely many terms exceed any
specified positive norm bound.  Thus this formula also specifies the
meaning of a series whose exponents have unbounded root denominators.

\begin{proposition}[Uniform relative narrowing]\label[proposition]{prop:kum-annulus}
The Kummer tower on $A_S[a,b]$ gives a natural complex
\begin{equation}\label{eq:kum-difference}
 \RGamma(A_S[a,b],\cO^\flat)
 \simeq
 \left[M_R[a,b]\xrightarrow{\gamma-1}M_R[a,b]\right],
 \qquad \gamma(T^u)=\epsilon^uT^u,
\end{equation}
in degrees $0,1$.  If
\begin{equation}\label{eq:kum-gap}
 a<a'<b'<b,\qquad
 a'/a\geq e^\kappa,\qquad b/b'\geq e^\kappa,
\end{equation}
restriction to $A_S[a',b']$ is chain homotopic to its
constant-monomial part.  It therefore factors through
$R[0]\oplus R[-1]$ in the derived category.  The homotopy is bounded,
linear over $R$, uniform for every compact profinite parameter, and
compatible with base change and a translation of the coordinate by
a relative section.
\end{proposition}

\begin{proof}
Adjoin all $p$-power roots of $T$.  The resulting perfectoid annulus
is a pro-\'etale $\Zp(1)$-torsor over $A_S[a,b]$.  Choose compatible
roots of the two scalar radii.  Tilting identifies its structural
ring with \cref{eq:kum-annular-algebra}, with the same radii: the
norm of a tilted coordinate equals the norm of its untilt.  The
$h$th term of the torsor nerve is the product of this perfectoid
annulus with $\Zp(1)^h$.  Its structural ring is the continuous
function ring on that compact parameter space.  All these objects
are affinoid perfectoid and are v-acyclic for $\cO^\flat$.
Descent consequently gives continuous group cochains.  The
two-term completed resolution for $\Zp$, justified for these
complete coefficients by \cref{lem:cts-comparison}, gives
\cref{eq:kum-difference}.  Notice that the topology here is the
Banach topology, not the discrete topology on the underlying ring.

For $u\ne0$, write $u=p^v m$, where $v\in\Z$ and $m$ is an integer
prime to $p$.  In characteristic $p$,
\[
 |\epsilon^u-1|=|\epsilon-1|^{p^v}.
\]
For $v\geq0$ this follows from Frobenius and from the unit linear
term of $(1+X)^m-1$; for $v<0$ take unique $p$-power roots.
Since $p^v\leq |u|_{\R}$, we obtain
\begin{equation}\label{eq:kum-denominator}
 |(\epsilon^u-1)^{-1}|
 \leq \exp(\kappa |u|_{\R}).
\end{equation}
The real absolute value of the exponent is essential in this
inequality.

Let $P_0$ denote projection onto the coefficient of $T^0$.  Define
the degree-$(-1)$ map from the source complex to the target by
\begin{equation}\label{eq:kum-homotopy}
 H\left(\sum_uc_uT^u\right)
   =\sum_{u\ne0}\frac{c_u}{\epsilon^u-1}T^u,
\end{equation}
on degree one, and set it equal to zero on degree zero.  For $u>0$
the ratio of the target and source Gauss weights is
$(b'/b)^u\leq e^{-\kappa u}$.  For $u<0$ it is
$(a'/a)^u\leq e^{-\kappa|u|}$.  These inequalities and
\cref{eq:kum-denominator} show that $H$ has operator norm at most
one.  In particular its series converges, including when root
denominators are unbounded.  Direct calculation gives
\[
 (\gamma-1)H+H(\gamma-1)=\res-P_0.
\]
This is the claimed chain homotopy.  Its entries are fixed scalars,
so it commutes with scalar extension and has the same bound for
$C(\mathcal D,R)$ coefficients.  A translation merely replaces
$T$ by the translated coordinate before applying the same formula.
\end{proof}

\begin{remark}\label{rem:kum-fixed-annulus}
The proposition concerns a map between two annuli.  It does not
assert that a fixed annulus has finite-dimensional cohomology.
Small denominators can contribute a large cokernel to
$\gamma-1$ at a fixed radius.  The homotopy disposes of these
classes only after the specified change of radii.
\end{remark}

\subsection{Discs, the line, and its boundary}

Write $D_S(c,r)=\{|T-c|\leq r\}$ for a disc about a relative
section.  Pullback from the base and restriction to the section
will be called the constant inclusion and evaluation, respectively.

\begin{proposition}[Relative line and disc calculations]\label[proposition]{prop:kum-line}
The following statements hold over $S$, naturally under base
change, with arbitrary retained compact profinite parameters.
\begin{enumerate}[label=\textup{(\roman*)}]
 \item There is a constant $A_p>1$, depending only on $p$ and the
 chosen Kummer generator, such that for $r_{\mathrm{out}}/r\geq A_p$ the map
 \[
 \RGamma(D_S(c,r_{\mathrm{out}}),\cO^\flat)
 \longrightarrow\RGamma(D_S(c,r),\cO^\flat)
 \]
 factors through evaluation $R[0]$ followed by the constant
 inclusion.  In particular,
 \[
 H^0(D_S(c,r),\cO^\flat)=R,
 \]
 and
 \[
 \colim_{r\to0}\RGamma(D_S(c,r),\cO^\flat)\simeq R[0].
 \]
 \item Ordinary and compactly supported cohomology of the line are
 \begin{equation}\label{eq:kum-line-values}
 \RGamma(\mathbb A_S,\cO^\flat)\simeq R[0],
 \qquad
 \RGamma_c(\mathbb A_S,\cO^\flat)\simeq R[-2].
 \end{equation}
 For two concentric support discs with a sufficiently large radius
 ratio, the map from the smaller support complex to the larger
 one has pure degree-two image, carried isomorphically into the
 final compact-support cohomology.  More precisely, if $A_r$ is
 the support complex for the disc of radius $r$ and
 $q_r:A_r\to R[-2]$ is extension of support to the whole line,
 the cone of $q_r$ maps cohomologically to zero after that fixed
 enlargement.  The required ratio is uniform, also under bounded
 translations of the centers.
 \item For $d$ ordered affine coordinates, the ordinary and section-germ
 complexes are $R[0]$, and the compact-support complex computed by
 cofinal boxes is $R[-2d]$. Maps between box-support bounds have pure
 degree-$2d$ image after a fixed enlargement in each coordinate.
 \item The analogous relative assertions with $\Fp$ coefficients
 hold with the base complex $\RGamma(S,\Fp)$ in place of $R$.
 For $S=\Spa(\Cflat,\cO_{\Cflat})\times\mathcal D$ this base
 complex is $C(\mathcal D,\Fp)[0]$.  With the chosen geometric
 orientations, extension of constants identifies these line,
 section-germ, and support computations with the corresponding
 $R_{\mathcal D}$ computations above.
\end{enumerate}
The degree-two support class is unchanged by a translation of its
section.  Under inclusion of finitely many disjoint support discs
into one larger disc, their degree-two classes add.
\end{proposition}

\begin{proof}
We first prove the restriction assertion.  By translation take
$c=0$.  Set
$\tau_p=|p|^{p/(p-1)}$.  Choose fixed constants $B>e^{2\kappa}$
and $A_p$ sufficiently large that, whenever $r_{\mathrm{out}}/r\geq A_p$, there
exists $s\in|C^\times|$ satisfying
\[
 r/s<\tau_p,\qquad Bs<r_{\mathrm{out}}.
\]
Increase $B$ slightly if needed to choose all radii strictly
separated.  Pick $\alpha\in C$ with $|\alpha|=s$.
Inside $D_S(0,r_{\mathrm{out}})$ there are annuli about $\alpha$, one contained
in the other, satisfying \cref{eq:kum-gap}, such that both contain
$D_S(0,r)$.  For example their inner and outer radii may be chosen
on opposite sides of $s$, each pair separated by a factor exceeding
$e^\kappa$, with outermost radius below $Bs$.  The restriction
between the discs factors through the restriction between these
annuli.  By \cref{prop:kum-annulus} this factors through
$R[0]\oplus R[-1]$.

The degree-one constant summand is the image of the mod-$p$
Kummer cocycle of $T-\alpha$, extended to $R$.  On $D_S(0,r)$,
\[
 T-\alpha=-\alpha(1-T/\alpha).
\]
Choose a $p$th root of $-\alpha$ in $C$.  The binomial series for
$(1-T/\alpha)^{1/p}$ converges because
$\|T/\alpha\|\leq r/s<\tau_p$; this convergence bound is
independent of the base.  Let $u(T)$ be this analytic $p$th root of
$T-\alpha$, and let $U$ denote the first root on the restricted
Kummer tower.  The ratio $U/u(T)$ is of the form $\zeta_p^{b}$
for a locally constant $\Fp$-valued function $b$ on that tower.
With the chosen generator, $\gamma(U)=\zeta_pU$, and hence
$(\gamma-1)b=1$.  Multiplication by this zero-cochain gives
an $R$-linear nullhomotopy of the constant degree-one class.
Thus the summand $R[-1]$ maps to zero in the derived category.
The remaining factorization through
$R[0]$ is evaluation: composing with the section of the smaller
disc shows that it agrees with evaluation on the larger disc.

For completeness, degree-zero sections on a fixed disc already
equal $R$.  Cover $D_S(0,r)$ by
\[
 \{a\leq |T|\leq r\},\qquad
 \{a\leq|T-\alpha|\leq r\},
 \quad 0<a<r,\quad |\alpha|=r.
\]
These annuli cover because the two omitted small discs are
disjoint.  In the complex \cref{eq:kum-difference} the kernel is
exactly $R$: every nonzero diagonal multiplier is an invertible
scalar.  The intersection contains a constant section
$T=\beta$ with $|\beta|=|\beta-\alpha|=r$, chosen using the
infinite residue field of $C$.  Hence two annular constants that
agree on the intersection are equal in $R$.  Sheaf descent proves
the assertion for the disc.  The restriction factorization and
exactness of filtered colimits now prove the shrinking-germ
statement in (i).

Choose an expanding sequence of closed discs covering the line,
with successive radius ratios at least $A_p$.  Sections on this
open exhaustion are the derived inverse limit of their section
complexes; using slightly enlarged open discs gives the same
limit.  By (i), its degree-zero cohomology system is the constant
system $R$, and every positive-degree system is pro-zero.
Equation~\eqref{eq:kum-milnor} proves the first assertion of
\cref{eq:kum-line-values}.

We next compute the germ of the exterior.  Let
$E_S(r)=\{|T|>r\}$.  The Kummer tower identifies its section
complex with
\begin{equation}\label{eq:kum-exterior}
 \left[M_R(r,\infty)\xrightarrow{\gamma-1}
 M_R(r,\infty)\right],
\end{equation}
where the coefficient ring consists of Laurent series converging
on every closed annulus contained in $r<|T|<\infty$.
To justify this assertion, exhaust $E_S(r)$ by a countable
increasing sequence of closed annuli.  Their coefficient
restriction maps have dense image, since finite Laurent
polynomials with bounded root denominator occur on every such
annulus and are dense on each target.  Apply
\cref{lem:kum-banach-limits} to the two terms of
\cref{eq:kum-difference}.  This gives exactly
\cref{eq:kum-exterior}, with its complete family of Gauss
seminorms.  The same argument works with compact parameters.

For $r'>e^\kappa r$, division by nonzero
$\epsilon^u-1$ defines a homotopy from
\cref{eq:kum-exterior} at $r$ to that at $r'$.
For negative exponents the increased inner radius absorbs
\cref{eq:kum-denominator}.  For positive exponents, to bound a
target seminorm with outer radius $b$, use the source seminorm
with outer radius larger than $e^\kappa b$; such a seminorm is
available because the outer radius is unbounded.  Thus the
homotopy is continuous and factors restriction through the
constant copies in degrees zero and one.  We conclude that
\begin{equation}\label{eq:kum-exterior-germ}
 \colim_{r\to\infty}\RGamma(E_S(r),\cO^\flat)
 \simeq R[0]\oplus R[-1].
\end{equation}
The map from ordinary sections of the line is the identity on
the first summand.  Its fiber is $R[-2]$, proving (ii).
Moreover the homotopy preserves the constant inclusion from
ordinary sections.  Taking its fiber proves the asserted pure
image statement for maps between support bounds, before taking
their final colimit.  Explicitly, the boundary of the constant
exterior cocycle gives a section $s_r:R[-2]\to A_r$ of $q_r$.
The fiber homotopy identifies a sufficiently separated transition
$A_r\to A_{r'}$ with $s_{r'}q_r$.  With the splitting
$A_r\simeq R[-2]\oplus\operatorname{fib}(q_r)$, the transition
is the identity on the first summand and zero on the second.
This proves the assertion about the cones as well.

Here is an explicit normalization of the support generator.
The constant degree-one Kummer cocycle of $T-c$ on an exterior
of $c$ maps, under the boundary map in
\cref{eq:kum-support-fiber}, to the section-support class.
For two bounded sections $c,c'$ and a sufficiently large exterior,
the quotient $(T-c')/(T-c)$ has a $p$th root: its difference from
$1$ has norm below $\tau_p$.  Their Kummer cocycles therefore
agree there.  This proves translation invariance and shows that
every section-support class maps to the same generator in
\cref{eq:kum-line-values}.  Excision for a finite disjoint union
of discs gives a direct sum of support complexes.  The map from
that direct sum to a common enlarged bound consequently sums
the section generators.

The arguments so far are relative: after a further affinoid
perfectoid base change all formulas are unchanged.  They
therefore identify the relative section, germ, and support
complexes as sheaves on the base v-site.  Apply them successively
in ordered coordinates.  The finite box-support diagram is the
iterated fiber of the restriction maps in those coordinates;
finite fibers commute, and the one-coordinate constant
inclusions and Kummer classes commute with every remaining
coordinate restriction.  The resulting complex is the tensor
product of $d$ copies of $R[-2]$, namely $R[-2d]$.
The same argument with evaluation gives ordinary and section-germ
complexes $R[0]$.  Taking a fixed sufficient radius ratio in each
coordinate gives the uniform pure-image statement in (iii).
This verifies the product assertion using the actual box fibers,
without assuming a compact-support comparison for arbitrary
analytic coefficient sheaves.

Finally use the absolute Artin--Schreier sequence on the v-site,
\begin{equation}\label{eq:kum-as}
 0\longrightarrow\Fp\longrightarrow\cO^\flat
 \xrightarrow{f\mapsto f^p-f}\cO^\flat\longrightarrow0.
\end{equation}
It is exact because its equation is solved by a finite \'etale
cover on every characteristic-$p$ perfectoid chart.  This is
also the sequence used in
\cite[proof of Theorem~5.1, p.~47]{scholzeHodge2013}.
The constant inclusions, evaluations, and chosen Kummer classes
are compatible with this sequence.  On the Laurent complexes
this compatibility is explicit: Frobenius sends $c_uT^u$ to
$c_u^pT^{pu}$, projection onto the constant term commutes with
it, and the homotopy does too, since
\[
 (\epsilon^u-1)^p=\epsilon^{pu}-1.
\]
Taking the fiber of Frobenius minus the identity in the preceding
relative computations therefore gives the corresponding base
complex $\RGamma(S,\Fp)$, with the same shifts.  Finite fibers
commute with the filtered colimits used here.

If $R=R_{\mathcal D}$, Frobenius minus the identity is surjective
on $R$.  Indeed the Artin--Schreier map on $\Cflat$ has a
continuous inverse branch on a sufficiently small ball about
each point of its target.  A finite clopen refinement of
$\mathcal D$ permits a continuous choice of these branches.
The kernel is $C(\mathcal D,\Fp)$, whose functions have finite
image.  This proves (iv).  For a general affinoid perfectoid base
we have kept the base complex, and have not asserted this global
surjectivity.  The finite-$p$ affine-line calculation also agrees
with the direct Kummer computation in
\cite[Example~2.19, proof of Lemma~2.20, Equation~(2.22)]
{achingerGillesNiziol2025}.
\end{proof}

\begin{remark}[Uniformity before a finite sum]\label[remark]{rem:kum-before-sum}
All support and restriction comparisons in
\cref{prop:kum-line} are made before any sum over translated
centers.  On finitely many relative charts, choose the largest
required radius ratio.  The same homotopies then apply with an
additional compact parameter and to every translated chart.
A map that is zero on a geometric cohomology row remains zero
after a finite sum of its translated maps.  Thus the estimates
are available before the finite coset grouping that produces
corestriction in \cref{sec:independent-tuples}.
\end{remark}

\subsection{Shrinking supports along a closed graph locus}

There are three different systems in use.  Shrinking
\emph{neighborhood germs} use restriction and a filtered colimit.
Enlarging compact supports uses maps from smaller support to
larger support and another filtered colimit.  Shrinking the
support bounds toward a closed locus uses those same support
maps as an inverse system.  The following result records the
last operation separately.

\begin{lemma}[Low degrees of shrinking graph supports]
\label[lemma]{lem:kum-low-support}
Let $X$ be a bounded relative chart in an untilted affine space over
a base v-sheaf $B$, with a slightly larger chart available for
collars, and let $Z\subset X$ be closed.  Assume the following
conditions on this fixed bound.
\begin{enumerate}[label=\textup{(\alph*)}]
 \item After a finite cover of a slightly enlarged base bound by
 qc charts, $Z$ has a cofinal decreasing sequence of
 neighborhoods $T_j$, each a finite disjoint union of relative
 discs, or of boxes in at least one normal affine coordinate.
 Their centers are actual analytic sections on those charts.
 \item On each chart, the clusters of $T_{j+1}$ refine those of
 $T_j$, and a fixed further refinement $j\mapsto j+N$ puts every
 child cluster in its parent with the radius ratio required in
 \cref{prop:kum-line}.  The same $N$ works for every $j$, for
 all compact parameters, and on the finitely many support
 shells used for excision.
 \item The cluster bounds have separated outer collars inside
 the enlarged output chart.  Their intersection is $Z$, and
 their complements admissibly exhaust $X\setminus Z$.
\end{enumerate}
For $\mathscr F=\cO^\flat$ or $\Fp$, put
\[
 E_j=\operatorname{fib}\bigl(\RGamma(X,\mathscr F)
              \longrightarrow\RGamma(X\setminus T_j,\mathscr F)
                       \bigr).
\]
The inverse systems $H^0(E_j)$ and $H^1(E_j)$ are pro-zero.
The same assertion holds relatively on the base and with compact
parameters.  Consequently the complex with support in $Z$ has
no cohomology in degrees zero or one, and deletion of $Z$
preserves relative degree-zero sections.  These conclusions hold
also after a bounded-below descent calculation for which the
finite chart refinements above can be made in each fixed total
degree.

For a closed locus $Z\subset Y$ in a possibly unbounded relative
untilted affine space, the same low-degree support vanishing and
degree-zero deletion statement hold if $Y$ has a compatible
countable admissible exhaustion by interleaved bounds $X_a$ on which
these hypotheses hold for $Z\cap X_a$.  The tube systems and their
required refinements may depend on $a$.
\end{lemma}

\begin{proof}
The map $E_{j+N}\to E_j$ goes from a smaller support bound to a
larger one.  Excision inside the separated collars identifies it
with a finite sum of maps from the support in child discs to
the support in their parent discs.  To see the excision used
here directly, cover an ambient chart by a neighborhood of its
support and the complement of a slightly smaller support bound.
The Mayer--Vietoris square identifies the two corresponding
fibers in \cref{eq:kum-support-fiber}.  Interleaving the inner
and outer bounds yields the stated collar excision.  It is
compatible with restriction on the base and with a further
finite cluster refinement.

By \cref{prop:kum-line}, after the fixed radius change each
child-to-parent map has pure geometric support image in degree
at least two.  The estimate is uniform before the finite sum,
by \cref{rem:kum-before-sum}.  Thus it induces zero on geometric
cohomology in degrees zero and one.  For $\Fp$ the base complex
in \cref{prop:kum-line}(iv) is connective, so its shift by a
support degree of at least two has the same low-degree
vanishing.  This proves the pro-zero assertion locally over
the base.

We spell out its passage through descent.  Use a qc covering
diagram in each cosimplicial degree.  Fix a total degree $i$.
Because the geometric complexes are bounded below by zero,
only cosimplicial degrees at most $i+1$ can affect the vanishing
in degrees up to $i$ and its comparison map.  Make the finitely
many required chart and collar refinements in these degrees
simultaneously.  The spectral sequence obtained by first taking
geometric cohomology has zero transition on the relevant
low-degree rows.  If necessary repeat the fixed refinement
once for each filtration piece.  A map zero on associated
graded pieces lowers the finite filtration, so this finite
composite is zero on the required total cohomology.  Hence the
total degree-zero and degree-one systems are pro-zero as
claimed.  This argument uses no common refinement for all
degrees of an unbounded descent diagram.

By (c), sections on $X\setminus Z$ are the derived inverse
limit of sections on the increasing complements
$X\setminus T_j$.  Finite fibers commute with this limit, so
\begin{equation}\label{eq:kum-closed-support-limit}
 \RGamma_Z(X,\mathscr F)\simeq\Rlim_j E_j.
\end{equation}
All $E_j$ are connective.  In
\cref{eq:kum-milnor}, the limit and first derived limit of the
pro-zero low-degree systems vanish.  This proves
$H^0\RGamma_Z=H^1\RGamma_Z=0$.  No surjectivity assertion about
the possibly large degree-two system is needed.  Finally the
long exact sequence for \cref{eq:kum-support-fiber}, with $Z$
in place of $K$, identifies degree-zero sections before and
after deletion.  The same proof applies to every retained
compact parameter, since all radius bounds and finite sums
were uniform in it.

For the last assertion use the interleaved open bounds of the
exhaustion and put
\[
 C_a=\RGamma_{Z\cap X_a}(X_a,\mathscr F).
\]
These support complexes have canonical restriction maps as $a$
varies, independently of the tube choices used to compute each one.
The fixed-bound argument gives $H^0(C_a)=H^1(C_a)=0$, and each $C_a$
is connective.  Both $X_a$ and $X_a\setminus Z$ admissibly exhaust
$Y$ and $Y\setminus Z$, respectively.  Taking their section limits
and commuting finite fibers with the derived inverse limit gives
\[
 \RGamma_Z(Y,\mathscr F)\simeq\Rlim_a C_a.
\]
The Milnor sequence again gives zero in degrees zero and one:
the relevant ordinary limits and $\varprojlim^1 H^0(C_a)$ are all
zero.  This uses the already formed support complexes $C_a$;
it requires no common tube depth for the different output bounds.
\end{proof}

\begin{remark}\label{rem:kum-tube-interface}
The fixed-bound hypotheses of \cref{lem:kum-low-support} will be verified in
\cref{prop:ind-tubes}. There a fixed lattice in the space of possible
relations indexes the clusters at every depth. Multiplication by $p$
rescales their centers and supplies the fixed rescaling in (b).
The lattice and its relation coordinates are defined in that argument.
The ordinary-function calculation then uses the separate exhaustion
assertion of the lemma to increase the output bound.
In particular, the lemma does not identify restriction of an arbitrary
v-sheaf to a closed analytic locus with germs of neighborhoods.
\end{remark}
